\chapter[Introducción]{Introducción\label{cap:introduccion}}

El estudio científico de las manchas solares no solo está vinculado con su importancia en la física solar, sino también con su valor práctico en tareas de monitoreo y registro observacional. La actividad fotosférica visible en estas estructuras aporta información sobre la evolución de regiones activas y, al mismo tiempo, ofrece contexto para interpretar condiciones de clima espacial capaces de afectar instrumentación sensible, sistemas electrónicos y entornos de observación científica. Sin embargo, en observatorios académicos y locales persiste un problema operativo: la adquisición de imágenes, el registro de metadatos, la identificación de grupos solares y la clasificación morfológica suelen ejecutarse como tareas separadas, con alta dependencia de revisión manual y con limitada trazabilidad entre la captura original y el resultado analítico. En ese marco, los avances relativos a herramientas que permitan adquirir, registrar y analizar observaciones de manera sistemática resultan relevantes no solo desde el punto de vista astronómico, sino también desde una perspectiva tecnológica y operativa \citep{hathaway2010solarcycle,usaf2003sunspot}.

La solución propuesta aborda este problema mediante un prototipo integrado de hardware, software e inteligencia artificial. En el plano instrumental, el sistema combina observación de luz blanca, cámara monocromática, atenuación óptica segura mediante una cuña de Herschel y una Raspberry Pi 5 como nodo local de adquisición y preprocesamiento. En el plano de software, incorpora servicios locales, sincronización de datos, almacenamiento estructurado y una plataforma web de asistencia. Sobre esta base se despliega una cadena de modelos de IA compuesta por detección de grupos solares, clasificación Zurich--McIntosh y módulos predictivos para el número de manchas solares y la distribución latitudinal de la actividad. Estos últimos no se incorporan solo como una salida complementaria, sino como una capa operacional de contexto: permiten ubicar cada observación dentro de una tendencia temporal más amplia, contrastar la actividad registrada localmente con patrones esperados del ciclo solar y aportar soporte interpretativo al usuario durante el seguimiento de la evolución fotosférica \citep{mcintosh1990classification,bornmann1994flarerates,pesnell2012predictions,clette2014revisiting}. De este modo, la tesis no se limita a entrenar modelos aislados, sino que propone una arquitectura completa para transformar una observación solar en un registro trazable, interpretable y contextualizado temporalmente.

Dicho esto, este trabajo se orienta al desarrollo de un prototipo funcional que integra obtención de imagen, registro y análisis asistido de manchas solares dentro de una misma arquitectura. El punto de partida de dicha integración son los dispositivos de instrumentación propios de la adquisición de imágenes heliográficas, como cámaras y filtros, compatibles con hardware de arquitectura embebida y con escenarios de observación local. Esta integración se apoya en la necesidad de vincular la calidad instrumental de la captura con flujos de procesamiento reproducibles, capaces de preservar trazabilidad entre la imagen, sus metadatos y las salidas analíticas posteriores \citep{scherrer2012hmi,bobra2014sharp,noaaSRS}.

La propuesta se apoya en una arquitectura de adquisición heliográfica y en una infraestructura de software capaz de sostener registro local, sincronización con base de datos y inferencia local en el dispositivo (\textit{on-device}) mediante visión por computadora. Sobre esa base se integran un detector basado en YOLO26n para el agrupamiento de manchas solares, una arquitectura jerárquica basada en ConvNeXtV2 Atto, complementada con etapas MLP, para inferir la clasificación Zurich--McIntosh y una etapa predictiva orientada al análisis temporal del número de manchas solares y de su distribución latitudinal. La etapa predictiva cumple una función operacional dentro del flujo: no sustituye la observación ni pretende producir pronósticos de clima espacial, pero permite relacionar las capturas locales con la dinámica de mayor escala del ciclo solar y ofrecer al usuario una lectura evolutiva de la actividad registrada \citep{podladchikova2012kalman,schmid2010dmd,kutz2016dmd}. De este modo, el trabajo articula observación, análisis asistido, persistencia de datos y soporte interpretativo dentro de una misma lógica operativa.

En este contexto, la tesis busca mostrar cómo decisiones de instrumentación, software, modelado, despliegue y análisis temporal pueden enlazarse de manera coherente dentro de un sistema de trabajo unificado \citep{clette2007wolf,hathaway2015solarcycle}.

\section{Objetivos}

\textbf{Objetivo general:}

Desarrollar un prototipo funcional para la deteccion y clasificacion local de manchas solares, integrando técnicas de visión por computadora, inteligencia artificial y análisis predictivo de actividad solar para observatorios locales.

\textbf{Objetivos específicos:}

\begin{itemize}
    \item Analizar registros historicos de manchas solares a partir de imágenes heliográficas y datos astronómicos.
    \item Construir un flujo reproducible de adquisición y preprocesamiento de imágenes que permita asociar observaciones solares con reportes NOAA SRS.
    \item Desarrollar un detector basado en YOLO26n para localizar grupos solares con latencia compatible con despliegue local.
    \item Diseñar un clasificador jerárquico para inferir los tres caracteres de la taxonomía Zurich--McIntosh preservando sus dependencias semánticas.
    \item Estimar parámetros físicos relevantes de las regiones activas, incluyendo posición heliográfica y evolución latitudinal.
    \item Implementar un módulo predictivo que aporte contexto operacional sobre la evolución temporal del índice de manchas solares y la distribución latitudinal de la actividad a partir de registros históricos suavizados.
    \item Integrar el sistema en una arquitectura híbrida con plataforma web y procesamiento inferencial local sobre dispositivos de acceso, incluyendo un nodo basado en Raspberry Pi 5 para operación en observatorio.
    \item Validar el funcionamiento integral del sistema mediante métricas de detección, clasificación, coherencia predictiva y consistencia física sobre datos históricos y escenarios de uso real.
\end{itemize}

\section{Alcance del trabajo}

El alcance del proyecto se enmarca en el desarrollo de un prototipo funcional orientado a observatorios locales en tierra, concebido para integrar en un mismo flujo operativo la captura de imágenes heliográficas, la persistencia de registros observacionales, el procesamiento inferencial basado en inteligencia artificial y la interpretación asistida de la actividad fotosférica \citep{hathaway2015solarcycle,sayez2023sunscc}. En ese sentido, el trabajo no se restringe al entrenamiento aislado de modelos, sino que aborda la articulación entre adquisición, almacenamiento, análisis y uso operativo dentro de una plataforma unificada. La intención central es demostrar que un sistema de estas características puede sostener tareas de observación local con criterios de trazabilidad, reproducibilidad y utilidad práctica para el estudio sistemático de manchas solares.

El prototipo contempla la integración de instrumentación compatible con entornos embebidos, la preparación de datos históricos y observacionales, el desarrollo de modelos de detección y clasificación morfológica, y la incorporación de modelos predictivos y herramientas de apoyo para el análisis temporal del número y la distribución de manchas solares \citep{ultralyticsYOLO26,deng2025mcintosh,pesnell2012predictions}. 

El alcance también incluye la dimensión operativa del prototipo. Por ello se consideran aspectos de despliegue local, interacción con una plataforma de software, persistencia de datos, ejecución inferencial sobre hardware de acceso relativamente bajo y generación de contexto predictivo para apoyar la interpretación de cada observación dentro de una tendencia temporal más amplia. Esta decisión resulta importante porque el valor del sistema no depende exclusivamente de la precisión de sus modelos, sino también de su capacidad para operar de forma coherente en escenarios de uso reales, donde intervienen restricciones de conectividad, disponibilidad computacional y continuidad de operación.

No forma parte del alcance la construcción de un observatorio solar completamente autónomo ni la sustitución de infraestructuras profesionales de monitoreo solar. Tampoco se persigue la formulación de un modelo físico exhaustivo de la dinámica fotosférica ni la elaboración de predicciones oficiales de clima espacial en sentido profesional; el componente predictivo se delimita, más bien, como una herramienta operacional interna para contextualizar la actividad observada dentro del flujo del prototipo. El aporte de esta tesis se concentra, en demostrar la viabilidad de una arquitectura funcional capaz de asistir tareas de registro, organización e interpretación local de observaciones solares mediante técnicas modernas de visión por computadora e inteligencia artificial. En consecuencia, el valor del trabajo reside en la integración de componentes instrumentales, computacionales y analíticos dentro de una propuesta unificada, concebida como base para desarrollos futuros de mayor complejidad.

\section{Estructura del documento}

La estructura del documento se presenta bajo una lógica de desarrollo progresivo: parte de los fundamentos físicos y observacionales del problema, avanza hacia las decisiones de diseño que condicionan la implementación, describe por separado las arquitecturas de hardware y software, y culmina con la validación experimental e integral del prototipo. Esta organización permite mantener una separación clara entre el contexto científico, la formulación metodológica, la construcción del sistema y la evaluación de sus resultados, incluyendo el papel del análisis predictivo como mecanismo de contextualización operacional y no únicamente como componente accesorio.

El Capítulo 2 desarrolla la revisión literaria que sustenta el trabajo. En su primera parte se presentan los conceptos necesarios para interpretar el fenómeno de estudio: actividad solar, regiones activas, manchas solares, modulación del ciclo solar y criterios observacionales empleados para describir grupos de manchas. Dentro de este marco se introduce la clasificación Zurich--McIntosh como referencia morfológica central, debido a que permite representar de forma estructurada la organización, penumbra y compactación de los grupos solares. En una segunda parte, el capítulo examina las condiciones de observación relevantes para el proyecto, incluyendo la diferencia entre observatorios terrestres y espaciales, las limitaciones introducidas por la atmósfera, y el papel de las cámaras astronómicas en la adquisición de imágenes de luz blanca. Finalmente, se revisa el estado del arte en detección, agrupamiento, clasificación y predicción de manchas solares, contrastando métodos clásicos basados en umbrales, morfología matemática y agrupamiento con enfoques recientes de aprendizaje profundo y análisis dinámico. Esta revisión establece la base técnica para justificar el uso de detectores YOLO, clasificadores jerárquicos y modelos predictivos basados en series temporales y distribución latitudinal \citep{ultralyticsYOLO26,sayez2023sunscc,podladchikova2012kalman}.

El Capítulo 3 presenta las consideraciones de diseño que orientan la propuesta. Su función es transformar los antecedentes científicos y tecnológicos en criterios concretos para construir el prototipo. En el plano instrumental, se analizan los requisitos de captura, atenuación óptica, compatibilidad con cámaras de astrofotografía y uso de una plataforma embebida como nodo local de adquisición. En el plano algorítmico, se define una estrategia modular para el análisis de imágenes solares: primero se localizan o agrupan regiones de interés mediante detección de objetos y luego se infiere la clasificación Zurich--McIntosh mediante un modelo jerárquico. En el plano predictivo, se formulan los criterios para representar la evolución temporal de la actividad mediante dos enfoques complementarios: un modelo escalar asociado al número de manchas solares y una representación espacio-temporal basada en la distribución latitudinal. De este modo, el capítulo funciona como puente entre la revisión literaria y las decisiones de implementación desarrolladas posteriormente, manteniendo la predicción como parte del flujo de interpretación y no como un resultado aislado \citep{schmid2010dmd,kutz2016dmd}.

El Capítulo 4 describe el diseño del hardware del sistema. Allí se documenta la arquitectura física que permite operar el prototipo en un entorno local de observación, con énfasis en la Raspberry Pi 5 como unidad de procesamiento embebido, los elementos de visualización de estado, el diseño mecánico de protección, la configuración de firmware, el sistema de captura de imágenes y la conectividad de red local. La descripción se orienta a mostrar cómo cada componente contribuye a un flujo de adquisición reproducible y seguro, desde la captura heliográfica hasta la disponibilidad de datos para su procesamiento posterior.

El Capítulo 5 desarrolla el diseño del software y del flujo computacional del prototipo. En primer lugar, describe la preparación de los datos de entrenamiento, la asociación de imágenes con reportes NOAA SRS, la construcción de conjuntos derivados y las limitaciones propias de la transferencia de dominio entre fuentes observacionales. Posteriormente, presenta la cadena de aprendizaje profundo empleada para el procesamiento de imágenes: un detector YOLO26n orientado al agrupamiento de manchas solares y un clasificador jerárquico basado en ConvNeXtV2 Atto, con etapas MLP para aprendizaje jerárquico, para inferir los tres componentes de la taxonomía Zurich--McIntosh. El capítulo incorpora además los módulos predictivos, incluyendo el modelo PVL con filtro de Kalman adaptativo y el modelo DMD espacio-temporal aplicado a una discretización latitudinal de la actividad. Finalmente, se describe la arquitectura de base de datos, la plataforma web, el despliegue en Next.js, la conversión de modelos a ONNX y la ejecución inferencial en el dispositivo del usuario. Esta organización permite presentar el software como una capa integradora entre datos, modelos, persistencia, interacción y contexto temporal de la actividad solar.

El Capítulo 6 reúne los resultados experimentales y la validación integral del sistema. La evaluación se organiza en dos niveles. El primero corresponde al desempeño de los modelos: detección con YOLO26n, clasificación jerárquica Zurich--McIntosh y comportamiento de los módulos predictivos asociados al número de manchas solares y a la distribución latitudinal, evaluados como parte del soporte operacional del sistema. El segundo nivel corresponde a la validación del prototipo, incluyendo pruebas del hardware, métricas de funcionamiento en Raspberry Pi 5 y evaluación de la plataforma web. Esta distinción permite valorar tanto la capacidad analítica de los modelos como la viabilidad práctica de integrarlos en un flujo real de observación, almacenamiento y consulta. El capítulo también explicita las limitaciones de la validación, especialmente en relación con el tamaño y naturaleza de los datos, la transferencia entre dominios observacionales, las condiciones de captura y el alcance de las predicciones.

El Capítulo 7 presenta la discusión, las conclusiones y las líneas de trabajo futuro. La discusión interpreta los resultados a partir del objetivo general de la tesis y analiza el grado en que la arquitectura propuesta logra integrar adquisición, registro, inferencia y soporte interpretativo. Las conclusiones sintetizan los aportes principales del trabajo: el diseño de un flujo local de observación y análisis, la incorporación de modelos de visión por computadora para detección y clasificación de manchas solares, la organización persistente de los registros y la inclusión de herramientas predictivas para contextualizar la evolución de la actividad fotosférica. Por último, los trabajos futuros identifican posibles mejoras en robustez observacional, ampliación de datos, validación con nuevas campañas de captura, optimización de inferencia en dispositivos de borde e integración con fuentes astronómicas externas.

Los anexos complementan el desarrollo principal mediante material técnico y visual de apoyo. Su inclusión permite preservar la continuidad argumental de los capítulos y, al mismo tiempo, documentar elementos relevantes para la implementación y comprensión del sistema. En conjunto, la estructura del documento busca demostrar de manera ordenada que el prototipo propuesto constituye una arquitectura funcional para transformar observaciones solares locales en registros trazables, procesables e interpretables mediante la convergencia de instrumentación heliográfica, software web, modelos de inteligencia artificial y análisis temporal de la actividad solar \citep{hathaway2010solarcycle,clette2014revisiting}.